% Final submission-oriented bibliography subset.
% 15 research papers total = 4 Chinese + 11 English, all published in 2022--2026.
% The broader research library remains in references.tex / REFERENCE_GUIDE.md.
\begingroup
\footnotesize
\setlength{\parskip}{0pt}
\begin{thebibliography}{99}
\setlength{\itemsep}{0.12ex}
\setlength{\parsep}{0pt}
\setlength{\topsep}{0.2ex}

% ===== 中文文献：算电协同、时空柔性、低碳与多元灵活性 =====
\bibitem{chen_idc_2022}
陈敏, 高赐威, 郭庆来, 李祖毅. 互联网数据中心负荷时空可转移特性建模与协同优化：驱动力与研究架构[J]. 中国电机工程学报, 2022, 42(19): 6945--6957.

\bibitem{zhang_flex_dro_2024}
张锞, 王旭, 杨宏坤, 蒋传文. 数据中心集群灵活边界下电力系统分布鲁棒优化调度方法[J]. 电力系统自动化, 2024, 48(7): 235--247.

\bibitem{wen_vpp_2024}
文淅宇, 朱继忠, 李盛林, 等. 基于时空协同的多数据中心虚拟电厂低碳经济调度策略[J]. 电力系统自动化, 2024, 48(18): 56--65.

\bibitem{zhou_multiflex_2024}
周世博, 周明, 孙黎滢, 等. 激发多元灵活性的数据中心协同优化运行方法[J]. 电网技术, 2024, 48(11): 4417--4426.

% ===== 总览与 Q1：地理分布式调度、工作负载统计与预测层级 =====
\bibitem{wu_survey_2025}
Wu Y, Tang S, Yu C, Yang B, Sun C, Xiao J, Wu H, Feng J. Task Scheduling in Geo-Distributed Computing: A Survey[J]. IEEE Transactions on Parallel and Distributed Systems, 2025, 36(10): 2073--2088.

\bibitem{cohen_dac_2022}
Cohen M C, Zhang R, Jiao K. Data Aggregation and Demand Prediction[J]. Operations Research, 2022, 70(5): 2597--2618.

\bibitem{zou_poissonity_2022}
Zou Q, Zhu Y, Tan Y, Chen W. Diagnosing the coexistence of Poissonity and self-similarity in memory workloads[J]. Journal of Network and Computer Applications, 2022, 205: 103455.

% ===== Q2：地理分布式任务时空调度与碳感知 =====
\bibitem{yang_repta_2026}
Yang L, Shahidehpour M, Chen X, Wang C, Fang X, Yang Q. High-speed REPTA algorithm for performance cost optimization in data centers considering workload delay tolerances[J]. Applied Energy, 2026, 404: 127156.

\bibitem{hanafy_carbonscaler_2023}
Hanafy W A, Liang Q, Bashir N, Irwin D, Shenoy P. CarbonScaler: Leveraging Cloud Workload Elasticity for Optimizing Carbon-Efficiency[J]. Proceedings of the ACM on Measurement and Analysis of Computing Systems, 2023, 7(3): 57:1--57:28.

\bibitem{sukprasert_limitations_2024}
Sukprasert T, Souza A, Bashir N, Irwin D, Shenoy P. On the Limitations of Carbon-Aware Temporal and Spatial Workload Shifting in the Cloud[C]//Proceedings of the 19th European Conference on Computer Systems. 2024: 924--941.

\bibitem{choudhury_mast_2024}
Choudhury A, Wang Y, Pelkonen T, et al. MAST: Global Scheduling of ML Training across Geo-Distributed Datacenters at Hyperscale[C]//18th USENIX Symposium on Operating Systems Design and Implementation. 2024: 563--580.

% ===== Q3：数据中心储能调度与物理可行性 =====
\bibitem{zhang_dc_bess_2025}
Zhang Y, Tang H, Li H, Wang S. Unlocking the flexibilities of data centers for smart grid services: Optimal dispatch and design of energy storage systems under progressive loading[J]. Energy, 2025, 316: 134511.

\bibitem{wang_exact_relax_2024}
Wang Q, Wu W, Lin C, Xu S, Wang S, Tian J. An exact relaxation method for complementarity constraints of energy storages in power grid optimization problems[J]. Applied Energy, 2024, 371: 123592.

% ===== Q4：算力—能源联合管理与分解求解 =====
\bibitem{guo_integrated_2025}
Guo H, Yu H, Wang M, Liu C, Li C. Integrated management of workloads and energy system for data centers[J]. Energy, 2025, 327: 136400.

\bibitem{wang_bcg_2026}
Wang C, Liu R, Pan E, Jin D, Fan X. Combining Benders decomposition and column generation for server allocation and scheduling in multiple queues[J]. Omega, 2026, 142: 103536.

\end{thebibliography}
\endgroup
